\documentclass[letterpaper,11pt,oneside]{article}

\usepackage{amsmath,amssymb,amsthm,mathtools}
\usepackage[final,colorlinks=true,linkcolor=blue,citecolor=red]{hyperref}
\usepackage{tikz}
\usepackage{booktabs}
\usepackage[DIV=14]{typearea}
\usepackage{caption}

\allowdisplaybreaks

\newtheorem{proposition}{Proposition}
\newtheorem{remark}[proposition]{Remark}

\newcommand{\ssp}{\hspace{1pt}}
\newcommand{\AZ}{\mathcal{A}}

\title{Shuffling and RSK at $q=0$ produce different conditional step distributions}
\author{Verification note}
\date{March 2026}

\begin{document}
\maketitle

%% =====================
%% TikZ style definitions
%% =====================
\tikzset{
  domN/.style={fill=yellow!75, draw=black!70, line width=0.4pt},
  domS/.style={fill=green!45!black!15, draw=black!70, line width=0.4pt},
  domE/.style={fill=blue!35, draw=black!70, line width=0.4pt},
  domW/.style={fill=red!35, draw=black!70, line width=0.4pt},
  badblock/.style={draw=red!80!black, line width=1.8pt, dashed, rounded corners=1pt},
  emptyblock/.style={draw=orange!80!black, line width=1.2pt, dotted, rounded corners=1pt,
    fill=orange!8},
  diamond boundary/.style={draw=black!50, line width=0.6pt, fill=black!4},
}

%% Macros: draw a horizontal domino at (u,v) with style
\newcommand{\hdom}[3]{\fill[#3] (#1,#2) rectangle ++(2,1);}
%% Vertical domino
\newcommand{\vdom}[3]{\fill[#3] (#1,#2) rectangle ++(1,2);}
%% A_2 diamond boundary path (TikZ coords with shift +2)
\newcommand{\aztecTwoBoundary}{
  (1,0) -- (3,0) -- (3,1) -- (4,1) -- (4,3) -- (3,3) -- (3,4) -- (1,4) --
  (1,3) -- (0,3) -- (0,1) -- (1,1) -- cycle
}
%% A_3 diamond boundary path (TikZ coords with shift +3)
\newcommand{\aztecThreeBoundary}{
  (2,0) -- (4,0) -- (4,1) -- (5,1) -- (5,2) -- (6,2) -- (6,4) -- (5,4) --
  (5,5) -- (4,5) -- (4,6) -- (2,6) -- (2,5) -- (1,5) -- (1,4) -- (0,4) --
  (0,2) -- (1,2) -- (1,1) -- (2,1) -- cycle
}

\section{Setup}

We compare two algorithms that produce uniform random domino tilings of the
Aztec diamond $\AZ_n$:

\smallskip

\noindent\textbf{EKLP Shuffling.}
Builds tilings iteratively via $\AZ_0 \to \AZ_1 \to \cdots \to \AZ_n$.
Each step $\AZ_k \to \AZ_{k+1}$ consists of three phases:
\begin{enumerate}
  \item \emph{Delete:} Find and remove all ``bad blocks'' --- colliding
    domino pairs (N-above-S or E-left-of-W in the same $2\times 2$ block).
  \item \emph{Slide:} Move each remaining domino one unit in its natural
    direction (N$\uparrow$, S$\downarrow$, E$\to$, W$\leftarrow$).
  \item \emph{Fill:} Each empty $2\times 2$ block in $\AZ_{k+1}$ is
    independently filled with a horizontal or vertical pair, probability $\tfrac12$ each.
\end{enumerate}

\noindent\textbf{RSK Growth Diagram at $q=0$.}
Fills an $(n{+}1)\times(n{+}1)$ staircase grid of partitions.
Each cell $(i,j)$ receives an independent fair Bernoulli bit, and the
partition at $(i,j)$ is determined by the VH bijection (which is deterministic
at $q=0$). The boundary of the sub-staircase at level $k$ encodes a tiling of
$\AZ_k$; the boundary layer from level $k$ to $k{+}1$ consists of exactly
$k{+}1$ cells, hence $k{+}1$ random bits.

\medskip

Both algorithms produce the correct uniform distribution on domino tilings.
The question is whether their \emph{conditional} distributions for individual
growth steps $\AZ_k \to \AZ_{k+1}$ coincide.

\section{Observation: the conditional distributions differ}

\begin{proposition}\label{prop:different}
  The EKLP shuffling and the RSK growth diagram at $q=0$ produce different
  conditional distributions $\Pr(\text{tiling of } \AZ_{k+1} \mid
  \text{tiling of } \AZ_k)$ for $k \ge 2$.
\end{proposition}

The mechanism is simple: the RSK step always uses exactly $k{+}1$ random
bits, while the shuffling step uses $k{+}1{+}m$ bits, where $m$ is the number
of bad blocks in the $\AZ_k$ tiling. When $m > 0$, the shuffling has strictly
more randomness in the conditional step.

\section{Example: $\AZ_2 \to \AZ_3$}

The Aztec diamond $\AZ_2$ has $2^{2 \cdot 3/2} = 8$ tilings.
Of these, 6 have no bad blocks ($m=0$) and 2 have exactly one bad block
($m=1$).

\subsection{The eight tilings of $\AZ_2$}

We display all eight tilings below.
Colors encode domino types: \tikz[baseline=-2pt]{\fill[domN] (0,0) rectangle (0.35,0.2);} N (north),
\tikz[baseline=-2pt]{\fill[domS] (0,0) rectangle (0.35,0.2);} S (south),
\tikz[baseline=-2pt]{\fill[domE] (0,0) rectangle (0.2,0.35);} E (east),
\tikz[baseline=-2pt]{\fill[domW] (0,0) rectangle (0.2,0.35);} W (west).
Bad blocks are marked with \tikz[baseline=-2pt]{\draw[badblock] (0,0) rectangle (0.4,0.3);}\,.

\bigskip

\begin{center}
\begin{tikzpicture}[scale=0.42]
%% === Tiling 1: all vertical EW ===
\begin{scope}[shift={(0,0)}]
  \fill[diamond boundary] \aztecTwoBoundary;
  \vdom{0}{1}{domW}  \vdom{3}{1}{domE}
  \vdom{1}{0}{domW}  \vdom{2}{0}{domE}
  \vdom{1}{2}{domW}  \vdom{2}{2}{domE}
  \node[below] at (2,-0.3) {\footnotesize $m{=}0$};
\end{scope}
%% === Tiling 2 ===
\begin{scope}[shift={(6,0)}]
  \fill[diamond boundary] \aztecTwoBoundary;
  \vdom{0}{1}{domW}  \vdom{3}{1}{domE}
  \hdom{1}{0}{domS}  \hdom{1}{1}{domN}
  \vdom{1}{2}{domW}  \vdom{2}{2}{domE}
  \node[below] at (2,-0.3) {\footnotesize $m{=}0$};
\end{scope}
%% === Tiling 3 ===
\begin{scope}[shift={(12,0)}]
  \fill[diamond boundary] \aztecTwoBoundary;
  \hdom{1}{0}{domS}  \hdom{1}{3}{domN}
  \vdom{0}{1}{domW}  \vdom{1}{1}{domE}
  \hdom{2}{1}{domS}  \hdom{2}{2}{domN}
  \node[below] at (2,-0.3) {\footnotesize $m{=}0$};
\end{scope}
%% === Tiling 4 ===
\begin{scope}[shift={(18,0)}]
  \fill[diamond boundary] \aztecTwoBoundary;
  \hdom{1}{0}{domS}  \hdom{1}{3}{domN}
  \hdom{0}{1}{domS}  \hdom{0}{2}{domN}
  \vdom{2}{1}{domW}  \vdom{3}{1}{domE}
  \node[below] at (2,-0.3) {\footnotesize $m{=}0$};
\end{scope}
\end{tikzpicture}

\bigskip

\begin{tikzpicture}[scale=0.42]
%% === Tiling 5: all horizontal ===
\begin{scope}[shift={(0,0)}]
  \fill[diamond boundary] \aztecTwoBoundary;
  \hdom{1}{0}{domS}  \hdom{1}{3}{domN}
  \hdom{0}{1}{domS}  \hdom{0}{2}{domN}
  \hdom{2}{1}{domS}  \hdom{2}{2}{domN}
  \node[below] at (2,-0.3) {\footnotesize $m{=}0$};
\end{scope}
%% === Tiling 6 ===
\begin{scope}[shift={(6,0)}]
  \fill[diamond boundary] \aztecTwoBoundary;
  \vdom{0}{1}{domW}  \vdom{3}{1}{domE}
  \vdom{1}{0}{domW}  \vdom{2}{0}{domE}
  \hdom{1}{2}{domS}  \hdom{1}{3}{domN}
  \node[below] at (2,-0.3) {\footnotesize $m{=}0$};
\end{scope}
%% === Tiling 7: BAD E-W block ===
\begin{scope}[shift={(12,0)}]
  \fill[diamond boundary] \aztecTwoBoundary;
  \hdom{1}{0}{domS}  \hdom{1}{3}{domN}
  \vdom{0}{1}{domW}  \vdom{1}{1}{domE}
  \vdom{2}{1}{domW}  \vdom{3}{1}{domE}
  \draw[badblock] (1.15,1.15) rectangle (2.85,2.85);
  \node[below] at (2,-0.3) {\footnotesize $m{=}1$ (E-W)};
\end{scope}
%% === Tiling 8: BAD N-S block ===
\begin{scope}[shift={(18,0)}]
  \fill[diamond boundary] \aztecTwoBoundary;
  \vdom{0}{1}{domW}  \vdom{3}{1}{domE}
  \hdom{1}{0}{domS}  \hdom{1}{1}{domN}
  \hdom{1}{2}{domS}  \hdom{1}{3}{domN}
  \draw[badblock] (1.15,1.15) rectangle (2.85,2.85);
  \node[below] at (2,-0.3) {\footnotesize $m{=}1$ (N-S)};
\end{scope}
\end{tikzpicture}
\end{center}

\subsection{The shuffling step for a bad-block tiling}

Consider Tiling~7 (E-W bad block).  The shuffling step
$\AZ_2 \to \AZ_3$ proceeds as follows:

\bigskip

\begin{center}
\begin{tikzpicture}[scale=0.38]
%% Panel 1: Original tiling with bad block
\begin{scope}[shift={(0,0)}]
  \fill[diamond boundary] \aztecTwoBoundary;
  \hdom{1}{0}{domS}  \hdom{1}{3}{domN}
  \vdom{0}{1}{domW}  \vdom{1}{1}{domE}
  \vdom{2}{1}{domW}  \vdom{3}{1}{domE}
  \draw[badblock] (1.15,1.15) rectangle (2.85,2.85);
  \node[below, font=\footnotesize\bfseries] at (2,-0.5) {(a) $\AZ_2$ tiling};
  \node[below, font=\scriptsize] at (2,-1.1) {1 bad block};
\end{scope}

%% Panel 2: After delete
\begin{scope}[shift={(6.5,0)}]
  \fill[diamond boundary] \aztecTwoBoundary;
  \hdom{1}{0}{domS}  \hdom{1}{3}{domN}
  \vdom{0}{1}{domW}  \vdom{3}{1}{domE}
  % empty region
  \draw[emptyblock] (1.1,1.1) rectangle (2.9,2.9);
  \node[below, font=\footnotesize\bfseries] at (2,-0.5) {(b) Delete};
  \node[below, font=\scriptsize] at (2,-1.1) {bad pair removed};
\end{scope}

%% Panel 3: After slide (in A_3 coords, shift +3)
\begin{scope}[shift={(14,0)}]
  \fill[diamond boundary, scale=0.75, shift={(0.83,0.83)}] \aztecThreeBoundary;
  % Slid dominoes (scaled)
  \begin{scope}[scale=0.75, shift={(0.83,0.83)}]
    \hdom{2}{0}{domS}   % S slid down
    \hdom{2}{5}{domN}    % N slid up
    \vdom{0}{2}{domW}    % W slid left
    \vdom{5}{2}{domE}    % E slid right
    % Empty region: 4x4 block from (1,1) to (5,5)
    \draw[emptyblock] (1.1,1.1) rectangle (2.9,2.9);
    \draw[emptyblock] (1.1,3.1) rectangle (2.9,4.9);
    \draw[emptyblock] (3.1,1.1) rectangle (4.9,2.9);
    \draw[emptyblock] (3.1,3.1) rectangle (4.9,4.9);
  \end{scope}
  \node[below, font=\footnotesize\bfseries] at (3.1,-0.5) {(c) Slide};
  \node[below, font=\scriptsize] at (3.1,-1.1) {into $\AZ_3$ frame};
\end{scope}

%% Panel 4: Fill annotation
\begin{scope}[shift={(22,0)}]
  \fill[diamond boundary, scale=0.75, shift={(0.83,0.83)}] \aztecThreeBoundary;
  \begin{scope}[scale=0.75, shift={(0.83,0.83)}]
    \hdom{2}{0}{domS}
    \hdom{2}{5}{domN}
    \vdom{0}{2}{domW}
    \vdom{5}{2}{domE}
    % Show fill choices as dice icons
    \draw[emptyblock] (1.1,1.1) rectangle (2.9,2.9);
    \draw[emptyblock] (1.1,3.1) rectangle (2.9,4.9);
    \draw[emptyblock] (3.1,1.1) rectangle (4.9,2.9);
    \draw[emptyblock] (3.1,3.1) rectangle (4.9,4.9);
    % coin flip symbols
    \node[font=\scriptsize\bfseries, orange!70!black] at (2,2) {?};
    \node[font=\scriptsize\bfseries, orange!70!black] at (2,4) {?};
    \node[font=\scriptsize\bfseries, orange!70!black] at (4,2) {?};
    \node[font=\scriptsize\bfseries, orange!70!black] at (4,4) {?};
  \end{scope}
  \node[below, font=\footnotesize\bfseries] at (3.1,-0.5) {(d) Fill};
  \node[below, font=\scriptsize] at (3.1,-1.1) {\textbf{4 blocks} $\to 2^4{=}16$};
\end{scope}
\end{tikzpicture}
\end{center}

\medskip

\noindent
\textbf{Comparison.}  For a tiling with $m = 0$ bad blocks, both algorithms
use $k{+}1 = 3$ random bits and produce $2^3 = 8$ equally likely extensions.
For a tiling with $m = 1$ bad block:

\begin{center}
\renewcommand{\arraystretch}{1.3}
\begin{tabular}{lcc}
  \toprule
  & \textbf{EKLP Shuffling} & \textbf{RSK Growth Diagram} \\
  \midrule
  Random bits per step & $k + 1 + m = 4$ & $k + 1 = 3$ \\
  Possible extensions  & $2^4 = 16$ & $2^3 = 8$ \\
  Probability per ext. & $\tfrac{1}{16}$ each & $\tfrac{1}{8}$ each \\
  \bottomrule
\end{tabular}
\end{center}

\smallskip

In the shuffling, deleting a bad block frees a $2\times 2$ region that later
gets re-filled with an independent coin flip, introducing one extra random
bit per bad block.  Since the RSK boundary layer has no analog of deletion,
its conditional distribution uses strictly fewer random choices.

\begin{remark}
  Despite the different conditional step distributions, both algorithms
  produce the same \emph{marginal} distribution at every level: uniform over
  all $2^{n(n+1)/2}$ tilings of $\AZ_n$.  The difference lies entirely in
  the coupling between consecutive levels.
\end{remark}

\section{Empirical verification}

We ran $5 \times 10^5$ samples per algorithm for the step $\AZ_2 \to \AZ_3$.
The partition-sequence encoding was used to identify tilings (with a
conjugation correction to match conventions between the two algorithms).

\smallskip

\noindent\textbf{Marginal distributions.}
Both algorithms produced all 8 tilings of $\AZ_2$ and all 64 tilings of
$\AZ_3$ with the correct uniform frequencies.  The total variation distance
between the marginal distributions was $< 0.002$ for $\AZ_2$ and $< 0.007$
for $\AZ_3$, consistent with sampling noise.  The chi-squared $p$-values
were $0.57$ and $0.44$, respectively.

\smallskip

\noindent\textbf{Conditional distributions.}
For the 6 states with $m = 0$: conditional TV distances were all below $0.01$
with chi-squared $p$-values $> 0.06$, consistent with identical distributions.

For the 2 states with $m = 1$: conditional TV distances were approximately
$\mathbf{0.50}$, with chi-squared $p$-values effectively zero ($\chi^2 >
40{,}000$).  The shuffling produced 16 distinct $\AZ_3$ extensions from each
of these states (each with probability $\approx 1/16$), while RSK produced
only~8 (each with probability $\approx 1/8$).

\smallskip

\noindent\textbf{Joint distribution.}
The joint distribution on $(\AZ_2, \AZ_3)$ pairs showed a TV distance of
$0.13$, with the shuffling producing 80 distinct pairs versus 64 for RSK.
The extra 16 pairs correspond to the 2 bad-block states $\times$ 8 extensions
that only the shuffling can reach.

\section{Discussion}

The EKLP shuffling and RSK growth diagram define two different
\emph{couplings} between $\AZ_k$ and $\AZ_{k+1}$ tilings that share the
same marginals but differ in their conditional structure.  The shuffling
coupling is ``more random'' due to the delete-and-refill mechanism: each bad
block contributes one extra fair coin flip to the transition.

For $k = 1$, neither algorithm produces bad blocks (the unique pair in
$\AZ_1$ is never a colliding pair), so the conditional distributions coincide.
Starting from $k = 2$, some $\AZ_k$ tilings contain bad blocks, and the
conditional distributions diverge.

The number of bad blocks $m$ depends on the specific $\AZ_k$ tiling.
For $\AZ_2$, the maximum is $m = 1$ (attained by 2 of the 8 tilings).
In general, $m$ can grow with $k$.

\end{document}
